%% ============================================================
%% 分片 35：原文第 12 节「Branes From Gauge Theory」（原文 glshort.tex 起于第 11186 行）
%% 由 tools/harvmac2tex.py 从 arXiv 源码（hep-th/0604151v3）自动生成。
%% 只译英文散文；公式、编号、\label、\cite 一律不改。
%% ============================================================

\section{Branes From Gauge Theory}
\label{sec:zapendix}

So far we have always considered branes in the effective
two-dimensional sigma-model with target $\MH$.  Our goal in this
concluding section is to show that the important branes actually
can be defined in four-dimensional gauge theory.

Branes of any interest always preserve at least one supersymmetry.
For example, they are $B$-branes in complex structure $J$ or
$A$-branes in complex structure $K$.  But the branes important for
geometric Langlands duality have special properties with respect
to all three complex structures. For example, the zerobrane on
$\MH(\LG,C)$ is a $(B,B,B)$-brane: it is a $B$-brane with respect
to complex structures $I,J,K$. Its mirror is a fiber of the
Hitchin fibration, which is complex with respect to $I$ and
Lagrangian with respect to $J$ and $K$. Thus it is a
$(B,A,A)$-brane. The canonical coisotropic brane is an
$(A,B,A)$-brane, and so is its mirror.  Branes of any of these
types preserve two topological supercharges (linear combinations
of which give, for example, for a brane of type $(B,B,B)$, the
requisite $B$-type supersymmetry in complex structures $I$, $J$,
$K$).

We will focus on these four kinds of brane and describe the
corresponding boundary conditions in the gauge theory.  The lift
to gauge theory is not necessarily unique; in fact, in some cases
we describe several distinct boundary conditions in the gauge
theory which upon reduction to two dimensions apparently become
equivalent, at least away from singularities.

To discuss the twisted theory as a topological field theory, one
replaces four-dimensional Minkowski space by a general
four-manifold $M$.  In topological field theory, it  is most
convenient to take $M$ to have Euclidean signature, for the
following reason. In going to a general four-manifold, in order to
preserve some supersymmetry, one ``twists'' by interpreting the
first four scalar fields $\phi_0,\dots,\phi_3$ as a section of the
tangent bundle of $M$. Since the scalar fields naturally have a
positive signature, this twisting is much more natural with
Euclidean signature on $M$. (Alternatively, one could possibly
make a Wick rotation in $\phi$ space, but this seems
unfelicitous.)  Additionally, any compact four-manifold admits a
positive signature metric, while admitting a metric of Lorentz
signature is a severe topological restriction.

For the present discussion, we will consider branes in a more
restricted situation.  As usual, we compactify to two dimensions
on a Riemann surface $C$.  This gives an effective two-dimensional
theory, and we want to understand how branes in this theory can be
interpreted in the underlying four-dimensional gauge theory.  Near
the boundary, there is a natural ``time'' direction (normal to $C$
and tangent to the boundary) and no major advantage in a Wick
rotation. On the contrary, we wish to avoid the Wick rotation in
order to make clear that all branes we consider are physically
sensible and unitary.

This being so, whenever a topological supercharge $Q$ is preserved
by our boundary conditions, its adjoint $Q^\dagger$ is also
preserved.  They obey a physical supersymmetry algebra
$\{Q,Q^\dagger\}=H$, with $H$ the Hamiltonian.  Since our boundary
conditions preserve two topological supercharges, they preserve a
total of four supersymmetries.  This is one-fourth of the 16
global supersymmetries of ${\EUN}=4$ super Yang-Mills theory, so
such branes are called 1/4 BPS branes.  From the point of view of
the sigma model, which has only one-half as much supersymmetry,
they are called 1/2 BPS branes.  Going back to gauge theory,
however, in some instances our boundary conditions in the gauge
theory preserve eight supersymmetries locally along $C$, and the
reduction to four supersymmetries comes from the twisting and
curvature of $C$.

To describe a brane in the two-dimensional effective theory, we
formulate this theory on (say) $\Bbb{R}^2_+$, with coordinates
$x^0$, $x^1$, where $-\infty< x^0< \infty$, $0\leq x^1 <\infty$. A
supersymmetric  brane is obtained by specifying a supersymmetric
boundary condition on $x^1=0$.  To obtain such a construction from
the underlying four-dimensional gauge theory, we simply consider
this  theory on $M=\Bbb{R}^2_+\times C$,  and again describe
supersymmetric boundary conditions at $x^1=0$.

We write $x^2,x^3$ for local coordinates on $C$.   As we take
$\Bbb{R}^2_+$ to be flat, the twisting only affects the two scalar
fields $\phi_2,\phi_3$ (or $A_6$, $A_7$; the notation was
explained in section 第~\ref{sec:reviewsym}~节). The supersymmetry left unbroken
by the compactification and twisting is generated by those spinors
$\epsilon$ that are covariantly constant on $C$ in the appropriate
sense.  The condition, in the twisted theory, is simply that
$\epsilon$ should obey % 式 (12.1)
\setcounter{equation}{0}
\begin{equation}\label{eq:ucuv}
\Gamma_{2367}\epsilon=\epsilon,
\end{equation}

as in eqn. \eqref{eq:ytytag}.

\subsection{General Properties Of Boundary Conditions}

First let us recall a few generalities about what a boundary
condition is supposed to be.

A boundary condition constrains the values of the fields (and
their normal derivatives, in the case of bosons that obey second
order equations of motion) in such a way that the boundary terms
in the equations of motion vanish.  For example, for gauge fields
on an $n$-manifold $M$ with metric $g$ of signature $-++\dots +$,
the minimal action is % 式 (12.2)
\setcounter{equation}{1}
\begin{equation}\label{eq:upco}
{{\cmmib I}}_B={1\over 2e^2}\int_M
d^nx \sqrt g\,\Tr \,F_{IJ}F^{IJ}.
\end{equation}
 $M$ is an $n$-manifold with
$x^0$ as the ``time'' direction and with a boundary defined by
$x^1=0$. The boundary terms in the variation $\delta {{\cmmib
I}}_B$ of the action are % 式 (12.3)
\setcounter{equation}{2}
\begin{equation}\label{eq:ofof}
(\delta {\cmmib
I}_B)_{bdry}={2\over e^2}\int_{\partial M}d^{n-1}x \sqrt
g\,\Tr\left(g^{IJ}\sum_{I,J\not=1}\delta A_I F_{1J}\right).
\end{equation}
  A
boundary condition must set to zero a linear combination of the
boundary values of $F_{I1}$ and $\delta A_I$, for $I\not= 1$, such
that $(\delta {{\cmmib I}}_B)_{bdry}=0$.  What it means to set
$\delta A_I=0$ for some $I$ is to specify the boundary values of
$A_I$ or in other words to impose Dirichlet boundary conditions on
$A_I$, with some prescribed boundary values. Conversely, in a
gauge with $A_1=0$ near the boundary of $M$, the condition that
$F_{I1}=0$ means that the normal derivative of $A_I$ vanishes.  So
$F_{I1}=0$ amounts to a gauge-invariant version of Neumann
boundary conditions on $A_I$. For quantization, we do not want to
overconstrain the boundary values, so we set to zero a minimal
linear combination of $\delta A_I$ and $F_{I1}$ such that $(\delta
{{\cmmib I}}_B)_{bdry}=0$.  (The boundary conditions therefore
define, in a suitable sense, a Lagrangian submanifold of the
boundary data.) It is also possible to add boundary terms to
${{\cmmib I}}_B$, thus modifying $(\delta {{\cmmib I}}_B)_{bdry}$;
the boundary conditions always define a maximal subspace of
allowed boundary values (of $A_I$ and $F_{I1}$, $I\not= 1$), such
that $(\delta {{\cmmib I}}_B)_{bdry}=0$. We call such a boundary
condition a hyperbolic boundary condition.

There is, of course, a similar story for fermions. The minimal
action is % 式 (12.4)
\setcounter{equation}{3}
\begin{equation}\label{eq:olpo}
{{\cmmib I}}_F=-{i\over e^2}\int_M d^nx \sqrt
g\,\Tr\,{\bar\lambda} \Gamma^I D_I\lambda,
\end{equation}
 to which one may add
additional boundary terms. Using the minimal action, the boundary
term in the variation of ${{\cmmib I}}_F$ is % 式 (12.5)
\setcounter{equation}{4}
\begin{equation}\label{eq:rolpo}
\delta
{{\cmmib I}}_F=-{i\over e^2}\int_{\partial M}d^{n-1}x\sqrt
g\,\Tr\,{\bar \lambda} \Gamma_1\delta\lambda.
\end{equation}
  For fermions, a
suitable boundary condition sets to zero the boundary values of
one-half the components of $\lambda$ (and therefore of
$\delta\lambda$) such that ${\bar \lambda}
\Gamma_1\delta\lambda=0$.

In the bulk theory, supersymmetry means (for $M$ flat or after
suitable twisting for $M$ curved) that there is a conserved
supercurrent $J_I$. For supersymmetric Yang-Mills theory,
% 式 (12.6)
\setcounter{equation}{5}
\begin{equation}\label{eq:noggin}
J_I={1\over 2}\G^{KL}F_{KL}\G_I\lambda.
\end{equation}
 In the
presence of a boundary, the condition that the boundary condition
is compatible with supersymmetry is that the normal component of
$J$ vanishes at the boundary, so that the flux of the supercurrent
does not disappear at the boundary.  More precisely, the condition
is % 式 (12.7)
\setcounter{equation}{6}
\begin{equation}\label{eq:hobbo}
\bar \epsilon J_1=0
\end{equation}
 for suitable supersymmetry
generators $\epsilon$.

This automatically ensures that the boundary conditions are
consistent with the supersymmetric variations of both bosons and
fermions.  If, therefore, one knows what boundary conditions one
wants on bosons and fermions, the condition $J_1=0$ is the basic
one.  In practice, since we do not know to begin with what
boundary conditions we want, we will start by postulating a simple
boundary condition on the fermions (chosen to ensure that ${\bar
\lambda}\Gamma_1\delta\lambda=0$), deduce the corresponding
boundary conditions on bosons to ensure that the boundary
conditions on the fermions are preserved by supersymmetry, and
then verify that $J_1=0$.

If $J_1=0$, we have found boundary conditions that will enable us
to define conserved supercharges (if $M$ is flat or in a suitably
twisted theory if $M$ is curved).  However, we want more.  A
boundary condition that forces supersymmetry to be spontaneously
broken will not be of much interest if our goal is to study
supersymmetric states or topological field theory (in which only
the supersymmetric states are of interest). Dirichlet boundary
conditions for bosons are generically incompatible with unbroken
supersymmetry. Dirichlet boundary conditions mean, as we have
discussed, that the boundary values of some components $A_I$ of
the gauge fields, and hence (possibly) of some components $F_{IJ}$
of the field strength, are specified.  We are mainly interested in
boundary conditions in which the boundary values of $F_{IJ}$ are
fixed in a way that is compatible with unbroken supersymmetry,
that is with the vanishing of the supersymmetry variation
% 式 (12.8)
\setcounter{equation}{7}
\begin{equation}\label{eq:undop}
\delta\lambda =\half\G^{IJ}F_{IJ}\epsilon
\end{equation}
 for suitable
supersymmetry generators $\epsilon$.

\subsection{Branes of Type $(B,B,B)$}

Our branes will always be constructed to preserve those
supersymmetries whose generators obey a suitable condition
$T\epsilon=\epsilon$ as well as\footnote{Replacing the condition
$T\epsilon=\epsilon$ with $T\epsilon=-\epsilon$ would lead to
nothing essentially new, as there always will be symmetries of the
theory that reverse the sign of $T$.  For example, in our first
example, $T=\Gamma_{04}$, such a symmetry would be  a reflection
in the $49$ plane.} $\Gamma_{2367}\epsilon=\epsilon$.
 These combined
conditions will allow four supersymmetries. $T$ will commute with
$\Gamma_{2367}$. In some examples, $T$ will be chosen so that
$T^2=1$.  Then the condition $T\epsilon=\epsilon$ would, by
itself, allow eight supersymmetries, and the boundary condition
will preserve those eight supersymmetries.  Even when that is so,
the compactification to two dimensions and the ``twisting'' will
impose the second condition $\Gamma_{2367}\epsilon=\epsilon$ on
unbroken supersymmetries, leaving only four supersymmetries in the
low energy theory.

We will explain illustrative choices of $T$ that lead to branes of
the various types.  We certainly do not claim to describe all
interesting constructions of branes in the four-dimensional gauge
theory, only some simple ones that are sufficient to exhibit in
the gauge theory the branes that will be most important in the
present paper.

To obtain branes of type $(B,B,B)$, we consider first $T=\G_{04}$
with the boundary condition % 式 (12.9)
\setcounter{equation}{8}
\begin{equation}\label{eq:unvu}
T\lambda\vert
=\lambda\vert.
\end{equation}
 For any field $\Phi$, we write $\Phi|$ for the
restriction of $\Phi$ to $x^1=0$. (However, we omit this symbol
when confusion seems unlikely.)  This boundary condition on
$\lambda$ is only compatible with supersymmetries such that the
supersymmetric transformation $\lambda\to \lambda+\delta\lambda$
preserves the condition $T\lambda\vert=\lambda\vert$.  The
requirement is that $(T-1)\delta\lambda=0$. Since
$\delta\lambda=\half\G^{IJ}F_{IJ}\epsilon$ and
$T\epsilon=\epsilon$, we have
 % 式 (12.10)
\setcounter{equation}{9}
\begin{equation}\label{eq:bogus}
(T-1)\delta\lambda =\half [T,\G^{IJ}F_{IJ}]\epsilon.
\end{equation}

Upon evaluating the commutator, we find that we require
% 式 (12.11)
\setcounter{equation}{10}
\begin{equation}\label{eq:nogus}
\sum_{I\not=0,4}(\Gamma_{0I}F_{4I}-\Gamma_{4I}F_{0I})\epsilon=0.
\end{equation}

Since $\epsilon$ is constrained precisely by
$\Gamma_0\Gamma_4\epsilon=\epsilon$, this condition is equivalent
to % 式 (12.12)
\setcounter{equation}{11}
\begin{equation}\label{eq:hogus}
F_{0I}+F_{4I}=0~~{\rm for}~I\not=0,4.
\end{equation}

To pick a hyperbolic boundary condition compatible with what has
just been described, we set (in a suitable gauge) the boundary
values of $A_0+A_4$ to zero, and for $A_I$, $I\not= 0,1,4$, we
impose Dirichlet boundary conditions with some specified,
time-independent, boundary values.  Finally, \eqref{eq:hogus}\ tells us that
$F_{01}+F_{41}=0$, corresponding to gauge-invariant Neumann
boundary conditions on $A_0+A_4$. (Since $A_0+A_4$ corresponds to
a ``null'' direction, it is possible in a hyperbolic boundary
condition to set to zero both $A_0+A_4$ and its normal derivative.
This depends on the minus sign in the signature, the fact that
$g^{IJ}={\rm diag}(-1,1,1,\dots,1)$ in \eqref{eq:ofof}.)

So far, starting with the boundary conditions on the fermions, we
have guessed what the boundary conditions on the bosons must be.
At this stage, we can verify the basic condition that the normal
component of the supercurrent vanishes on the boundary:
% 式 (12.13)
\setcounter{equation}{12}
\begin{equation}\label{eq:jolygo}
{\bar\epsilon}\G^{MN}F_{MN}\G_1\lambda=0.
\end{equation}
 Because
$\G_{04}\lambda=\lambda$ and $\G_{04}\epsilon=\eps$, the only
components of $F_{MN}$ that contribute are those with precisely
one of $M$ and $N$ equal to 0 or 4; using this, one can verify
that \eqref{eq:jolygo}\ is obeyed if $F_{0I}+F_{4I}=0$, $I\not= 0,4$.

This implies that the boundary conditions on both fermions and
bosons are compatible with supersymmetry. We have already seen
this for the fermions, and we can readily verify it for the
bosons.  The supersymmetry variation of the bosons
% 式 (12.14)
\setcounter{equation}{13}
\begin{equation}\label{eq:guffo}
\delta A_I=i\bar\epsilon\Gamma_I\lambda
\end{equation}
 implies, with
$(\Gamma_{04}-1)\lambda=(\G_{04}-1)\epsilon=0$, that $\delta
A_I=0$ for $I\not= 0,4$, and $\delta A_0+\delta A_4=0$, whence
$\delta(F_{0I}+F_{4I})=0$ for $I\not=0,4$.  Hence our condition
$F_{0I}+F_{4I}=0$, $I\not= 0,4$, is consistent with supersymmetry.

So far, we have obtained a boundary condition that ensures the
existence of conserved supercharges.  We also would like the
boundary condition to allow for existence of a supersymmetric
state.  For this, the Dirichlet boundary conditions on $A_I$,
$I\not=0,1,4$, are highly constrained.  Given the supersymmetric
variation $\delta \lambda= \half
\sum_{IJ}\Gamma^{IJ}F_{IJ}\epsilon$, we want to pick the boundary
values such that % 式 (12.15)
\setcounter{equation}{14}
\begin{equation}\label{eq:humbo}
\sum_{I,J\not=
0,1,4}\G^{IJ}F_{IJ}\epsilon=0.
\end{equation}
 If this were supposed to be so for
all $\epsilon$, we would have to set $F_{IJ}=0$, $I,J\not= 0,1,4$.
Since $\epsilon$ is constrained to
$\Gamma_{2367}\epsilon=T\epsilon=\epsilon$, the constraint on $F$
is less severe.  It suffices to take $F$ to be selfdual in the
subspace generated by $x^2,x^3,x^6$, and $x^7$. Thus, we must
impose the Hitchin equations, which in this notation read
% 式 (12.16)
\setcounter{equation}{15}
\begin{equation}\label{eq:ibob}
\begin{aligned}F_{23}-F_{67}& =0\\ F_{27}-F_{36}& = 0 \\
 F_{26}-F_{73}& = 0 . \\\end{aligned}
\end{equation}
  Thus, the boundary values of the gauge fields
define a point $p$ in $\MH$, the moduli space of Higgs bundles.
Once we impose \eqref{eq:ibob}, \eqref{eq:humbo}\ is obeyed if in addition $F_{IJ}|=0$
unless both $I$ and $J$ are in the set ${\EUR}=\{2,3,6,7\}$. In
particular, vanishing of $F_{IJ}|$ for $I\in {\EUR}$, $J\notin
{\EUR}$ implies that $A_J|$, for $J=5,8,9$ is covariantly constant
(on the Riemann surface $C$ with local coordinates $x^2,x^3$) and
commutes with the Higgs field, whose components are $\phi_2=A_6 $
and $\phi_3=A_7$.   If $p$ is a smooth point in $\MH$
corresponding to an irreducible Higgs bundle, these conditions
imply that $A_J|=0$ for $J=5,8,9$.  This furthermore implies the
vanishing of $F_{IJ}|$ for $I,J\notin {\EUR}$. So for an
irreducible Higgs bundle, the boundary condition (of this type) is
uniquely determined.  If $p$ corresponds to a reducible Higgs
bundle, we only learn in general that $A_5,$ $A_8$, and $A_9$
generate symmetries of the Higgs bundle and commute with each
other.

The boundary conditions we have obtained so far correspond to a
zerobrane on $\MH$ supported at the point $p$.  This is a basic
type of $(B,B,B)$-brane.  As expected, such a brane is unique if
$p$ corresponds to a smooth point in $\MH$.

\par\medskip{\it Another Construction Of The Zerobrane}

Somewhat surprisingly, there is another gauge theory construction
of zerobranes on $\MH$.  It seems to lead to the same brane in the
low energy theory as long as one is at a smooth point in $\MH$,
but may lead to something new in the case of a zerobrane supported
at a singular point in $\MH$.

Let us replace $T$ by $\tilde T = T
\Gamma_{2367}=\Gamma_{023467}$. The supersymmetries preserved by
the boundary condition will now be those whose generators obey
$\tilde T\epsilon=\epsilon$.  So the boundary condition in the
four-dimensional gauge theory will preserve a different set of
eight supercharges than what we had before. However, the unbroken
supersymmetries in the low energy theory are those that obey the
combined conditions $T\epsilon=\Gamma_{2367}\epsilon=\epsilon$.
These conditions are invariant under replacing $T$ by $\tilde T$,
so the change in boundary conditions will not change the unbroken
supersymmetry in the low energy theory and the construction will
give a new $(B,B,B)$-brane.

We will choose the boundary condition on the fermions to be
$\tilde T\lambda\vert = -\lambda\vert$. Asking for this condition
to be invariant under supersymmetry, we require that
$\{\Gamma^{IJ}F_{IJ},\tilde T\}\epsilon=0$. Let us assume in
addition that $\Gamma_{2367}\eps=\eps$. Then the supersymmetry
condition is satisfied if % 式 (12.17)
\setcounter{equation}{16}
\begin{equation}\label{eq:kipo}
\sum_{I,J\in
\{0,2,3,4,6,7\}}\Gamma^{IJ}F_{IJ}\vert\epsilon=0
\end{equation}
 and
% 式 (12.18)
\setcounter{equation}{17}
\begin{equation}\label{eq:kipok}
\sum_{I,J\in
\{1,5,8,9\}}\Gamma^{IJ}F_{IJ}\vert\epsilon=0.
\end{equation}
 \eqref{eq:kipo}\ is obeyed if
the $A_I$ obey Dirichlet boundary conditions with $A_4\vert=0,$
and with the boundary values of $A_I$, $I\in \{2,3,6,7\}$
constrained to obey the Hitchin equations \eqref{eq:ibob}.

Note that if $\eps$ does not obey $\Gamma_{2367}\eps=\eps$, the
supersymmetry constraint is not satisfied, in general. This means
that even in flat spacetime, this $(B,B,B)$-brane  preserves only
four supercharges.

To obey \eqref{eq:kipok},  we {\it cannot} ask for $F_{IJ}|=0$, $I,J\in
\{1,5,8,9\}$ as these conditions are overdetermined. Indeed,
$F_{15}|=F_{18}|=0$ means that $A_5$ and $A_8$ obey Neumann
boundary conditions and are free to fluctuate on the boundary;
this being so, we cannot also ask for $F_{58}\vert=0$. However,
the condition $\tilde T\epsilon=\epsilon$, together with the
chirality condition $\hat\Gamma\epsilon=\epsilon$ of
ten-dimensional super Yang-Mills theory (here
$\hat\Gamma=\Gamma_0\Gamma_1\cdots \Gamma_9$), implies
$\Gamma_{1589}\epsilon=-\epsilon$.  With this constraint on
$\epsilon$, \eqref{eq:kipok}\ is equivalent to an anti-selfduality condition
on the boundary values in the $1589$ plane:
% 式 (12.19)
\setcounter{equation}{18}
\begin{equation}\label{eq:ubb}
\begin{aligned}F_{15}\vert+F_{89}\vert & = 0 \\
                  F_{18}\vert+F_{95}\vert & = 0 \\
                  F_{19}\vert+F_{58}\vert & = 0. \\\end{aligned}
\end{equation}

This is a modified version of Neumann boundary conditions on
$A_5,$ $A_8$, and $A_9$, and is {\it not} overdetermined.

Since $A_2,A_3, A_6$, and $A_7$ obey Dirichlet boundary conditions
determined by a specific solution of the Hitchin equations, what
we obtain in this way is a zerobrane supported at a point $p$ in
$\MH$.  What distinguishes it from the previous construction of
the zerobrane is the boundary condition on $A_4,A_5,A_8$, and
$A_9$, as well as the fact that in flat spacetime the former
construction preserves half as many supersymmetries as the latter.
If $p$ is a smooth point in $\MH$ corresponding to an irreducible
Higgs bundle, the fields $A_4,A_5,A_8$, and $A_9$ have no zero
modes and vanish in the low energy theory, regardless of the
microscopic boundary conditions. Hence the new construction of the
zerobrane should be equivalent to the old one as long as $p$ is a
smooth point in $\MH$.  If $p$ corresponds to a reducible Higgs
bundle whose stabilizer is a subgroup of $G$ of positive rank,
then $A_4$, $A_5$, $A_8$, and $A_9$ do have zero modes -- which
generate symmetries of the Higgs bundle -- and there is no reason
to expect the two zerobranes to be equivalent.

One can get yet another $(B,B,B)$-brane by taking the boundary
condition for the fermions to be $\tilde T\lambda\vert =
\lambda\vert$. It turns out that it corresponds to a space-filling
brane with a flat gauge field. Since this brane does not play a
major role in the present paper, we omit the details.
